\FigureLayoutDeclare{img/2021/zhejiang/q7_fig1.png}{6.0cm}{9bp 10bp 0bp 0bp}
\FigureLayoutDeclare{img/2021/zhejiang/q11_fig1.png}{4.8cm}{10bp 10bp 10bp 10bp}
\FigureLayoutDeclare{img/2021/zhejiang/q16_fig1.png}{5.0cm}{10bp 10bp 10bp 11bp}

\chapter{2021年浙江卷}

\section{单选题}

\begin{problem}
设集合
\(A=\{x\mid x\ge 1\}\)，\(B=\{x\mid -1<x<2\}\)，则
\(A\cap B=\)（\quad）

\choices
    {\(\{x\mid x>-1\}\)}
    {\(\{x\mid x\ge 1\}\)}
    {\(\{x\mid -1<x<1\}\)}
    {\(\{x\mid 1\le x<2\}\)}
\end{problem}

\begin{problem}
已知 \(a\in\R\)，\((1+a\i)\i=3+\i\)（\(\i\) 为虚数单位），则
\(a=\)（\quad）

\choices
    {\(-1\)}
    {\(1\)}
    {\(-3\)}
    {\(3\)}
\end{problem}

\begin{problem}
已知非零向量 \(\bs{a}\)，\(\bs{b}\)，\(\bs{c}\)，则“\(\bs{a}\cdot\bs{c}=\bs{b}\cdot\bs{c}\)”是“\(\bs{a}=\bs{b}\)”的（\quad）

\choices
    {充分不必要条件}
    {必要不充分条件}
    {充分必要条件}
    {既不充分也不必要条件}
\end{problem}

\begin{problem}
某几何体的三视图如图所示（单位：\(\mathrm{cm}\)），则该几何体的体积（单位：\(\mathrm{cm}^3\)）是（\quad）

\bitmapfigure[max width=6.0cm]{img/2021/zhejiang/q4_fig1.png}

\choices
    {\(\frac{3}{2}\)}
    {\(3\)}
    {\(\frac{3\sqrt{2}}{2}\)}
    {\(3\sqrt{2}\)}
\end{problem}

\begin{problem}
若实数 \(x\)，\(y\) 满足约束条件 \(\begin{cases}
x+1\ge 0\\
x-y\le 0\\
2x+3y-1\le 0
\end{cases}\)，则 \(z=x-\frac{1}{2}y\) 的最小值是（\quad）

\choices
    {\(-2\)}
    {\(-\frac{3}{2}\)}
    {\(-\frac{1}{2}\)}
    {\(\frac{1}{10}\)}
\end{problem}

\begin{problem}
如图，已知正方体
\(ABCD-A_{1}B_{1}C_{1}D_{1}\)，
\(M\)，\(N\) 分别是 \(A_{1}D\)，\(D_{1}B\) 的中点，则（\quad）

\bitmapfigure[max width=6.0cm]{img/2021/zhejiang/q6_fig1.png}

\choices
    {直线 \(A_{1}D\) 与直线 \(D_{1}B\) 垂直，直线 \(MN\parallel\) 平面 \(ABCD\)}
    {直线 \(A_{1}D\) 与直线 \(D_{1}B\) 平行，直线 \(MN\perp\) 平面 \(BDD_{1}B_{1}\)}
    {直线 \(A_{1}D\) 与直线 \(D_{1}B\) 相交，直线 \(MN\parallel\) 平面 \(ABCD\)}
    {直线 \(A_{1}D\) 与直线 \(D_{1}B\) 异面，直线 \(MN\perp\) 平面 \(BDD_{1}B_{1}\)}
\end{problem}

\begin{problem}
已知函数
\(f(x)=x^{2}+\frac{1}{4}\)，\(g(x)=\sin x\)，则图象为如图的函数可能是（\quad）

\bitmapfigure[max width=6.0cm]{img/2021/zhejiang/q7_fig1.png}

\choices
    {\(y=f(x)+g(x)-\frac{1}{4}\)}
    {\(y=f(x)-g(x)-\frac{1}{4}\)}
    {\(y=f(x)g(x)\)}
    {\(y=\dfrac{g(x)}{f(x)}\)}
\end{problem}

\begin{problem}
已知 \(\alpha\)，\(\beta\)，\(\gamma\) 是互不相同的锐角，则在
\(\sin\alpha\cos\beta\)，\(\sin\beta\cos\gamma\)，\(\sin\gamma\cos\alpha\)
三个值中，大于 \(\frac{1}{2}\) 的个数的最大值是（\quad）

\choices
    {\(0\)}
    {\(1\)}
    {\(2\)}
    {\(3\)}
\end{problem}

\begin{problem}
已知 \(a\)，\(b\in\R\)，\(ab>0\)，函数
\(f(x)=ax^{2}+b\)（\(x\in\R\)）. 若
\(f(s-t)\)，\(f(s)\)，\(f(s+t)\) 成等比数列，则平面上点 \((s,t)\) 的轨迹是（\quad）

\choices
    {直线和圆}
    {直线和椭圆}
    {直线和双曲线}
    {直线和抛物线}
\end{problem}

\begin{problem}
已知数列 \(\{a_n\}\) 满足
\(a_1=1\)，\(a_{n+1}=\dfrac{a_n}{1+\sqrt{a_n}}\)（\(n\in\N^*\)），
记数列前 \(n\) 项和为 \(S_n\). 则（\quad）

\choices
    {\(\frac{1}{2}<S_{100}<3\)}
    {\(3<S_{100}<4\)}
    {\(4<S_{100}<\frac{9}{2}\)}
    {\(\frac{9}{2}<S_{100}<5\)}
\end{problem}

\section{填空题}

\begin{problem}
我国古代数学家赵爽用弦图给出了勾股定理的证明. 弦图是由四个全等的直角三角形和中间的一个小正方形拼成的大正方形（如图所示）.
若直角三角形直角边分别为 \(3\) 与 \(4\)，记大正方形面积为 \(S_1\)，小正方形面积为 \(S_2\)，则 \(\frac{S_1}{S_2}=\fillinblank{}\).

\bitmapfigure[max width=4.8cm]{img/2021/zhejiang/q11_fig1.png}
\end{problem}

\begin{problem}
已知 \(a\in\R\)，函数 \(f(x)=\begin{cases}
x^{2}-4,&x>2,\\
|x-3|+a,&x\le 2,
\end{cases}\) 若 \(f\left[f(\sqrt{6})\right]=3\)，则 \(a=\fillinblank{}\).
\end{problem}

\begin{problem}
已知
\((x-1)^3+(x+1)^4=x^4+a_1x^3+a_2x^2+a_3x+a_4\)，则
\(a_1=\fillinblank{}\)，
\(a_2+a_3+a_4=\fillinblank{}\).
\end{problem}

\begin{problem}
在 \(\triangle ABC\) 中，\(\angle B=60^\circ\)，\(AB=2\)，\(M\) 是 \(BC\) 中点，\(AM=2\sqrt{3}\)，则
\(AC=\fillinblank{}\)，\(\cos\angle MAC=\fillinblank{}\).

\bitmapfigure[max width=4.7cm]{img/2021/zhejiang/q14_fig1.png}
\end{problem}

\begin{problem}
袋中有 \(4\) 个红球，\(m\) 个黄球，\(n\) 个绿球，现从中任取两球，记取出的红球数为 \(\xi\).
若取出的两个球都是红球概率为 \(\frac{1}{6}\)，
一红一黄概率为 \(\frac{1}{3}\)，则
\(m-n=\fillinblank{}\)，\(E(\xi)=\fillinblank{}\).
\end{problem}

\begin{problem}
已知椭圆 \(\frac{x^2}{a^2}+\frac{y^2}{b^2}=1\)（\(a>b>0\)），焦点为
\(F_1(-c,0)\)，\(F_2(c,0)\)（\(c>0\)）.
若过 \(F_1\) 的直线和圆
\((x-\frac{c}{2})^2+y^2=c^2\) 相切，与椭圆在第一象限交于点 \(P\)，且 \(PF_2\perp x\) 轴，
则该直线斜率为 \(\fillinblank{}\)，椭圆的离心率为 \(\fillinblank{}\).

\bitmapfigure[max width=5cm]{img/2021/zhejiang/q16_fig1.png}
\end{problem}

\begin{problem}
已知平面向量 \(\bs{a}\)，\(\bs{b}\)，\(\bs{c}\)（\(\bs{c}\ne\bs{0}\)）
满足 \(|\bs{a}|=1\)，\(|\bs{b}|=2\)，\(\bs{a}\cdot\bs{b}=0\) 且 \((\bs{a}-\bs{b})\cdot\bs{c}=0\).
记向量 \(\bs{d}\) 在 \(\bs{a}\)，\(\bs{b}\) 方向上的投影分别为 \(x\)，\(y\)，
\(\bs{d}-\bs{a}\) 在 \(\bs{c}\) 方向上的投影为 \(z\)，则
\(x^2+y^2+z^2\) 的最小值为 \(\fillinblank{}\).
\end{problem}

\section{解答题}

\begin{problem}
设函数 \(f(x)=\sin x+\cos x\)（\(x\in\R\)）.

\begin{enumerate}
\item 求函数 \(y=\bigl[f(x+\tfrac{\pi}{2})\bigr]^2\) 的最小正周期；

\item 求函数 \(y=f(x)f(x-\tfrac{\pi}{4})\) 在 \([0,\frac{\pi}{2}]\) 上的最大值.
\end{enumerate}
\end{problem}

\begin{problem}
如图，在四棱锥 \(P-ABCD\) 中，底面 \(ABCD\) 为平行四边形，\(\angle ABC=120^\circ\)，
\(AB=1\)，\(BC=4\)，\(PA=\sqrt{15}\).
设 \(M\)，\(N\) 分别为 \(BC\)，\(PC\) 中点，\(PD\perp DC\)，\(PM\perp MD\).

\begin{enumerate}
\item 证明：\(AB\perp PM\)；

\item 求直线 \(AN\) 与平面 \(PDM\) 所成角的正弦值.
\end{enumerate}

\bitmapfigure[max width=6.0cm]{img/2021/zhejiang/q19_fig1.png}
\end{problem}

\begin{problem}
已知数列 \(\{a_n\}\) 前 \(n\) 项和为 \(S_n\)，\(a_1=-\frac{9}{4}\)，且
\(4S_{n+1}=3S_n-9\).

\begin{enumerate}
\item 求数列 \(\{a_n\}\) 的通项；

\item 设 \(\{b_n\}\) 满足
\(3b_n+(n-4)a_n=0\)（\(n\in\N^*\)），记其前 \(n\) 项和为 \(T_n\). 若
\(T_n\le\lambda b_n\) 对任意 \(n\in\N^*\) 恒成立，求 \(\lambda\) 的取值范围.
\end{enumerate}
\end{problem}

\begin{problem}
如图，已知 \(F\) 是抛物线 \(y^2=2px\)（\(p>0\)）的焦点，\(M\) 为该抛物线准线与 \(x\) 轴交点，且 \(|MF|=2\).

\begin{enumerate}
\item 求抛物线方程；

\item 设过 \(F\) 的直线交抛物线于 \(A\)，\(B\)，斜率为 \(2\) 的直线 \(l\) 与直线 \(MA\)，\(MB\)，\(AB\)，\(x\) 轴依次交于 \(P\)，\(Q\)，\(R\)，\(N\)，且 \(|RN|^2=|PN|\cdot|QN|\)，求直线 \(l\) 在 \(x\) 轴上的截距的范围.
\end{enumerate}

\bitmapfigure[max width=6.0cm]{img/2021/zhejiang/q21_fig1.png}
\end{problem}

\begin{problem}
设 \(a\)，\(b\in\R\)，且 \(a>1\)，函数 \(f(x)=a^x-bx+\e^2\)（\(x\in\R\)）.

\begin{enumerate}
\item 求函数 \(f(x)\) 的单调区间；

\item 若对任意 \(b>2\e^2\)，函数 \(f(x)\) 有两个不同零点，求 \(a\) 的取值范围；

\item 当 \(a=\e\) 时，证明：对任意 \(b>\e^4\)，函数 \(f(x)\) 有两个不同零点 \(x_1\)，\(x_2\)，使 \(x_2>\frac{b\ln b}{2\e^2}x_1+\frac{\e^2}{b}\).
\end{enumerate}
\end{problem}

